\section{Results: fragment recombination builds the library}
\label{sec:res:fragments}

\plan{\textbf{Claim.} Designs from the trained policy can be cut at fixed
boundaries and recombined across parents, so that $\sum_i |F_i|$ synthesized DNA
parts yield $\prod_i |F_i|$ testable variants --- and the recombinants keep
enough of the parental pass rate that the resulting library is worth building.

\textbf{Status: \statusblocked{}.} This is the one section of the paper with no
data behind it yet. Every other claim in the document is either measured or a
rewrite of something measured. Written out in full here so the experiments it
needs are unambiguous.

\textbf{Why it is the load-bearing section.} Section~\ref{sec:res:frontier}
buys a better \emph{distribution}; on its own that is a sampling result. This
section is what converts a distribution into something you can physically order,
and it is the only reason the PETase libraries in
Section~\ref{sec:res:petase} can be the size they are. Without it the paper is a
benchmark result with an application attached.}

\subsection*{Block 1 --- The combinatorial argument, stated precisely}

\begin{planlist}
  \item Set up the notation: $L$ residues split into $k$ fragments at fixed
        boundaries $[s_i, e_i]$; $F_i$ the set of distinct sequences observed at
        fragment $i$ across the policy's output; a recombinant is one choice per
        slot, joined.
  \item The arithmetic: $\sum_i |F_i|$ parts, $\prod_i |F_i|$ recombinants. With
        $k = 5$ and $|F_i| = 40$ that is 200 parts and $1.0 \times 10^8$
        variants. \textbf{State the actual numbers used, not this example.}
  \item Golden Gate assembly is what makes the product physically reachable:
        one pot, orthogonal overhangs at the boundaries, no per-variant
        synthesis. Methods carries the adapter design.
  \item \textbf{The assumption the whole section rests on, named up front:}
        fragments are treated as if they contribute independently to
        foldability. They do not, exactly --- the whole point of an inverse
        folding model is that residues are coupled. Block 3 measures how badly
        the approximation fails, and that measurement is the result.
\end{planlist}
\todo{fragments --- state the actual $k$, $|F_i|$ and library size used}

\subsection*{Block 2 --- Where the boundaries go}

\begin{planlist}
  \item Boundaries are chosen on the structure, not the sequence: cut in loops
        and between secondary-structure elements, never through the active site.
  \item \textbf{Pre-register the comparison:} boundaries placed on
        secondary-structure breaks should retain more pass rate than boundaries
        placed at evenly spaced sequence positions. If they do not, the
        structural rationale is decorative and should be dropped from the paper
        rather than defended.
  \item Report the fixed-residue set --- catalytic residues held constant across
        all fragments, so every recombinant retains the active site by
        construction.
\end{planlist}
\todo{fragments --- boundary placement rule + the structural-vs-uniform control}

\subsection*{Block 3 --- Do recombinants fold? (the experiment that must be run)}

\plan{\textbf{This is the measurement the section exists for.} Everything above
is setup and everything below is consequence. Design, in the order it should be
run:

\begin{enumerate}\setlength{\itemsep}{0pt}
  \item Sample $n$ designs from a trained \cleo{} policy with no selection
        applied, on an AME backbone where the parental pass rate is already
        known and high (M0097 is the obvious choice --- it is the backbone every
        other measurement in the paper uses).
  \item Split each at the fragment boundaries. Keep only fragments drawn from
        \emph{passing} parents, since that is what one would actually order.
  \item Draw recombinants uniformly from the product space, excluding any
        recombinant that reproduces a parent exactly.
  \item Fold best-of-5 and score on the identical code path as
        Section~\ref{sec:res:frontier}. \textbf{Same pipeline, same reference,
        or the comparison to the parental rate is not a comparison.}
\end{enumerate}

\textbf{The endpoint is retention:} recombinant pass rate divided by parental
pass rate. Everything else is secondary.

\textbf{Pre-registered predictions, written before the run:}
\begin{itemize}\setlength{\itemsep}{0pt}
  \item Retention is well above the rate for recombining across \emph{random}
        MPNN samples at matched diversity. If it is not, fragmentation is
        adding nothing that temperature would not.
  \item Retention falls monotonically with the number of slots swapped away
        from a single parent. A recombinant differing at one slot should behave
        nearly like its parent; one differing at all $k$ should be worst. A flat
        curve would mean the fragments are effectively independent, which would
        be a stronger result than we expect and should be reported as such.
  \item Recombinants reach mutational combinations absent from the parental set
        --- measured as $U$ over the recombinant library against $U$ over the
        parents at matched design count.
\end{itemize}

\textbf{What kills the claim:} retention below the point where the library
is filterable in practice. If the parental rate is \SI{57}{\percent} and
recombinants pass at \SI{2}{\percent}, the product space is nominal rather than
real, and the honest paper reports the fragment library as a screening funnel
with a stated hit rate instead of as a $10^8$-member library. Decide the
threshold now, not after seeing the number.}

\todo{fragments --- RUN THIS. Selection-free sample $\rightarrow$ split $\rightarrow$
  recombine $\rightarrow$ fold best-of-5 $\rightarrow$ retention curve vs.\ slots
  swapped}

\subsection*{Block 4 --- Relation to manufacturing-aware generative models}

\plan{\citet{weinstein2026manufacturing} is the nearest prior art and the
section must engage it directly rather than in passing. Their move is to make
the generative model manufacturing-aware from the outset --- factorized to match
what combinatorial DNA chemistry can physically sample --- and they reach
${\sim}10^{16}$ designs for antibodies, T-cell antigens and Taq polymerase.

\textbf{The honest framing.} They get there first and at larger scale. What is
different is the regime, and it is a real difference rather than a hedge:

\begin{itemize}\setlength{\itemsep}{0pt}
  \item \textbf{Their objective comes from a natural distribution.} 300 million
        observed human antibodies define what a good sample looks like, so the
        model can be constrained to the chemistry from the start without giving
        anything up. For a de novo enzyme there is no such distribution --- the
        backbone has no natural homologs --- so the objective has to be
        constructed, which is what Section~\ref{sec:res:frontier} does.
  \item \textbf{Order of operations.} They factorize first and train the
        factorized model; we train an unfactorized policy for catalytic geometry
        and factorize afterwards. That makes retention (Block 3) our problem and
        not theirs --- their samples are manufacturable by construction, ours
        have to be shown to survive the cut.
  \item \textbf{Scale is not the axis we win on and the section should say so
        plainly.} $10^{16}$ is out of reach here and it does not need to be
        reached: a $10^{8}$ library that is enriched for catalytic geometry is a
        different object from a $10^{16}$ library that is not.
\end{itemize}

\textbf{Do not write this block as a competitive comparison.} The two results
compose --- a manufacturing-aware architecture trained against a catalytic
reward is the obvious next thing, and saying so is more useful than claiming a
win.}

\todo{fragments --- write Block 4 once Block 3 lands; needs the library-size
  and retention numbers to be concrete about the $10^{8}$ vs $10^{16}$ point}
